\begin{table}[htbp]
\centering
\caption{Do Colorado Border Counties Opt In to Retail Marijuana at a Higher Rate?}
\label{tab:politicaleconomy}
\small
\begin{tabular}{lccc}
\toprule
& Interior counties & Border counties & \\
& ($n=50$) & ($n=14$) & \\
\midrule
Share ever with $\geq$1 dispensary & 72.0\% & 42.9\% & \\
\midrule
Specification & Coefficient & $p$-value & $n$ \\
\midrule
No controls & $-0.291$ & $0.051$ & 64 \\
+ log(population) & $-0.263^{**}$ & $0.036$ & 64 \\
+ log(population), pop.\ $\geq$5{,}000 only & $-0.258$ & $0.105$ & 52 \\
\bottomrule
\end{tabular}
\begin{minipage}{0.95\textwidth}
\vspace{0.3em}
\footnotesize
\textit{Notes:} Outcome is a binary indicator for whether a Colorado county ever had an operating
dispensary between 2015 and 2025, the best available proxy for effective local exclusion of retail
marijuana under Amendment~64's opt-out authority (no compiled municipal-ordinance dataset exists; this
proxy conflates active local bans with no operator ever entering a low-demand market). Sample is the full
64-county Colorado universe, not only the 42 counties in the paper's revenue panel. Linear probability
model, HC1 standard errors. \textbf{The sign is opposite to the fiscal-capture prediction}: border
counties are less, not more, likely to have ever permitted retail marijuana. $^{**}p<0.05$.
\end{minipage}
\end{table}
